% ─── Chapter 6: Conclusion ───────────────────────────────────────────────────

% Slide 1: Summary of Achievements & Methodology
\begin{frame}{Summary of Achievements \& Implemented Methods}
  \small
  \begin{columns}[T]
    \begin{column}{0.48\textwidth}
      \begin{block}{Key Achievements}
        \footnotesize
        \begin{itemize}
          \setlength{\itemsep}{3pt}
          \item \textbf{High-Fidelity Solver}: Built a 2D Isogeometric Analysis (IGA) physical solver ($N=6{,}564$ DOFs).
          \item \textbf{Hyper-Reduction}: Deployed ECSW on only 28 cells, yielding a \textbf{$64.5\times$ speedup} with $0.054\%$ error.
          \item \textbf{Digital Twin}: Interactive client-side JS application achieving real-time solve times ($< 1.5\text{ ms}$).
        \end{itemize}
      \end{block}
    \end{column}

    \begin{column}{0.48\textwidth}
      \begin{block}{Implemented Methods \& References}
        \footnotesize
        \begin{itemize}
          \setlength{\itemsep}{2pt}
          \item \textbf{POD / SVD}: Dimension reduction basis.
          \newline \textit{Ref: Benner et al. (2015)}
          \item \textbf{ECSW Hyper-Reduction}: Sparse mesh quadrature. \textit{Ref: Farhat et al. (2014)}
          \item \textbf{Diffeomorphic Pullback}: Shape-parameter mapping preserving mesh topology.
          \item \textbf{Newmark-$\beta$}: Transient dynamic solver.
        \end{itemize}
      \end{block}
    \end{column}
  \end{columns}
\end{frame}

% Slide 2: Unexplored Alternatives & Future Work
\begin{frame}{Alternatives \& Future Perspectives}
  \small
  \begin{columns}[T]
    \begin{column}{0.48\textwidth}
      \begin{block}{Other Methods (Unexplored)}
        \footnotesize
        \begin{itemize}
          \setlength{\itemsep}{2pt}
          \item \textbf{DEIM / MDEIM}: Discrete Empirical Interpolation for nonlinear terms.
          \newline \textit{Refs: Chaturantabut (2010), Bonomi (2017)}
          \item \textbf{Gappy POD}: Regression-based state reconstruction.
          \item \textbf{Machine Learning}: Autoencoder-based nonlinear ROMs or PINNs.
          \item \textbf{Multi-Field ROMs}: Damage-plasticity decompositions. \textit{Ref: Kehls (2025)}
        \end{itemize}
      \end{block}
    \end{column}

    \begin{column}{0.48\textwidth}
      \begin{block}{Future Research Paths}
        \footnotesize
        \begin{itemize}
          \setlength{\itemsep}{2pt}
          \item \textbf{Nonlinear Kinematics}: Large displacements using Green-Lagrange strain directly in reduced space.
          \item \textbf{3D NURBS}: Trivariate solid formulations and multi-patch assemblies.
          \item \textbf{Adaptive Bases}: Online local basis update/enrichment (e.g., Kriging).
        \end{itemize}
      \end{block}
    \end{column}
  \end{columns}
\end{frame}
